% Añadir paquetes
\usepackage{amsmath, amsfonts}  % Para las expresiones matemáticas
\usepackage{xcolor}             % Para utilizar colores en la escritura
\usepackage{lipsum}             % Para generar texto aleatorio
\usepackage{acronym}            % Para utilizar acrónimos
\usepackage{graphicx}           % Para insertar imágenes en el documento
\usepackage{float}              % Para posicionar imágenes en el documento
\usepackage{multirow}           % Para juntar múltiples lineas en una tabla
\usepackage{appendix}           % Para el diseño de los anexos
\usepackage{blindtext}          % Para generar texto aleatorio
\usepackage[ruled]{algorithm2e} % Para generar pseudocódigo
\usepackage{pdfpages}           % Para  añadir pdfs
\usepackage{booktabs}
\usepackage{longtable}          % Para tablas que ocupan varias páginas
\usepackage{listings}           % Para bloques de código con lstlisting
\usepackage{titlesec}
\usepackage{caption}
\usepackage{tabularx}
\usepackage{hyperref}
\usepackage[section]{placeins}
\setcounter{secnumdepth}{4}

\titleformat{\paragraph}
{\normalfont\normalsize\bfseries}{\theparagraph}{1em}{}
\titlespacing*{\paragraph}
{0pt}{3.25ex plus 1ex minus .2ex}{1.5ex plus .2ex}

\lstdefinestyle{mystyle}{
    basicstyle=\ttfamily\small,
    breaklines=true,
    frame=single,
    framesep=6pt,
    framexleftmargin=4pt,
    framexrightmargin=4pt,
    framextopmargin=4pt,
    framexbottommargin=0pt,
    numbers=left,
    numberstyle=\tiny,
    aboveskip=0.75\baselineskip,
    belowskip=0.25\baselineskip,
    showstringspaces=false,
    tabsize=4
}

% Reducir el espacio caption -> next paragraph
\setlength{\belowcaptionskip}{-10pt}  
\captionsetup[table]{aboveskip= 12pt}

% Copiado de mathrsfs.sty para generar slash zero
\DeclareSymbolFont{rsfs}{U}{rsfs}{m}{n}
\DeclareSymbolFontAlphabet{\mathscrsfs}{rsfs}
\usepackage{stackengine}[2013-09-11]
\newcommand\slashzero{\stackinset{c}{}{c}{}{/}{0}}


\renewcommand{\figurename}{Figura}   % Cambiar el nombre a la figura
\renewcommand{\tablename}{Tabla}     % Cambiar el nombre a la tabla
\renewcommand{\chaptername}{Capítulo}  % Cambiar el nombre a la tabla
\SetAlgorithmName{Algoritmo}{algoritmo}{Lista de Algoritmos} % Cambiar el nombre al algoritmo

% Generación del color del texto "Memoria"
\usepackage[dvipsnames]{xcolor}
\definecolor{myblue}{rgb}{0.282, 0.239, 0.545}


% Generación del subsubsection{}
\makeatletter
\setcounter{secnumdepth}{3}
\titlespacing{\subsubsection}{0em}{.5em}{0em}%[0pt]
\ifct@cthesis@hangsubsection
        \titleformat{\subsubsection}[hang]%
            {\usekomafont{subsubsection}}%
            {\color{ctcolorblack}\thesubsubsection\hspace*{10pt}}%
            {0pt}%
            {\raggedright}%
            [\phantomsection]
\else
        \titleformat{\subsubsection}[block]%
            {\usekomafont{subsubsection}}%
            {\color{ctcolorblack}\thesubsubsection\hspace*{10pt}}%
            {0pt}%
            {\raggedright}%
            [\phantomsection]
\fi
\makeatother